\documentclass[10pt,a4paper]{amsart}
\usepackage{amsmath,amssymb,amsthm,graphicx}
% For commutative diagrams we use Xy-pic
\usepackage[all]{xy}

% We use 2cell for 2-commutative diagrams.
\xyoption{2cell}
% Theorem environments.
%
\theoremstyle{plain}
\newtheorem{theorem}[subsection]{Theorem}
\newtheorem{proposition}[subsection]{Proposition}
\newtheorem{lemma}[subsection]{Lemma}

\theoremstyle{definition}
\newtheorem{definition}[subsection]{Definition}
\newtheorem{example}[subsection]{Example}
\newtheorem{exercise}[subsection]{Exercise}
\newtheorem{situation}[subsection]{Situation}

\theoremstyle{remark}
\newtheorem{remark}[subsection]{Remark}
\newtheorem{remarks}[subsection]{Remarks}

\numberwithin{equation}{subsection}

% Macros
%
\def\lim{\mathop{\mathrm{lim}}\nolimits}
\def\colim{\mathop{\mathrm{colim}}\nolimits}
\def\Spec{\mathop{\mathrm{Spec}}}
\def\Hom{\mathop{\mathrm{Hom}}\nolimits}
\def\Ext{\mathop{\mathrm{Ext}}\nolimits}
\def\SheafHom{\mathop{\mathcal{H}\!\mathit{om}}\nolimits}
\def\SheafExt{\mathop{\mathcal{E}\!\mathit{xt}}\nolimits}
\def\Sch{\mathit{Sch}}
\def\Mor{\mathop{\mathrm{Mor}}\nolimits}
\def\Ob{\mathop{\mathrm{Ob}}\nolimits}
\def\Sh{\mathop{\mathit{Sh}}\nolimits}
\def\NL{\mathop{N\!L}\nolimits}
\def\CH{\mathop{\mathrm{CH}}\nolimits}
\def\proetale{{pro\text{-}\acute{e}tale}}
\def\etale{{\acute{e}tale}}
\def\QCoh{\mathit{QCoh}}
\def\Ker{\mathop{\mathrm{Ker}}}
\def\Im{\mathop{\mathrm{Im}}}
\def\Coker{\mathop{\mathrm{Coker}}}
\def\Coim{\mathop{\mathrm{Coim}}}

% Boxtimes
%
\DeclareMathSymbol{\boxtimes}{\mathbin}{AMSa}{"02}

%
% Macros for moduli stacks/spaces
%
\def\QCohstack{\mathcal{QC}\!\mathit{oh}}
\def\Cohstack{\mathcal{C}\!\mathit{oh}}
\def\Spacesstack{\mathcal{S}\!\mathit{paces}}
\def\Quotfunctor{\mathrm{Quot}}
\def\Hilbfunctor{\mathrm{Hilb}}
\def\Curvesstack{\mathcal{C}\!\mathit{urves}}
\def\Polarizedstack{\mathcal{P}\!\mathit{olarized}}
\def\Complexesstack{\mathcal{C}\!\mathit{omplexes}}
% \Pic is the operator that assigns to X its picard group, usage \Pic(X)
% \Picardstack_{X/B} denotes the Picard stack of X over B
% \Picardfunctor_{X/B} denotes the Picard functor of X over B
\def\Pic{\mathop{\mathrm{Pic}}\nolimits}
\def\Picardstack{\mathcal{P}\!\mathit{ic}}
\def\Picardfunctor{\mathrm{Pic}}
\def\Deformationcategory{\mathcal{D}\!\mathit{ef}}
\usepackage{cleveref}
\crefname{originallemma}{original lemma}{original lemmas}
\Crefname{originallemma}{Original lemma}{Original lemmas}
\makeatletter
\newlabel{coherent-lemma-quasi-coherence-higher-direct-images}{{01XJ}{source}{}{sourcepointer-1}{}}
\newlabel{coherent-lemma-quasi-coherence-higher-direct-images@cref}{{[originallemma][0][]01XJ}{[1][0][]source}{}{}{}}
\makeatother
\crefname{originallemma}{original lemma}{original lemmas}
\Crefname{originallemma}{Original lemma}{Original lemmas}
\makeatletter
\newlabel{descent-lemma-higher-direct-images-small-etale}{{071N}{source}{}{sourcepointer-2}{}}
\newlabel{descent-lemma-higher-direct-images-small-etale@cref}{{[originallemma][0][]071N}{[1][0][]source}{}{}{}}
\makeatother
\crefname{originalreference}{original reference}{original references}
\Crefname{originalreference}{Original reference}{Original references}
\makeatletter
\newlabel{lemma-cohomology-with-support-triangle}{{0A45}{source}{}{sourcepointer-3}{}}
\newlabel{lemma-cohomology-with-support-triangle@cref}{{[originalreference][0][]0A45}{[1][0][]source}{}{}{}}
\makeatother
\crefname{originalreference}{original reference}{original references}
\Crefname{originalreference}{Original reference}{Original references}
\makeatletter
\newlabel{proposition-quasi-coherent-sheaf-fpqc}{{03OG}{source}{}{sourcepointer-4}{}}
\newlabel{proposition-quasi-coherent-sheaf-fpqc@cref}{{[originalreference][0][]03OG}{[1][0][]source}{}{}{}}
\makeatother
\crefname{originalreference}{original reference}{original references}
\Crefname{originalreference}{Original reference}{Original references}
\makeatletter
\newlabel{theorem-zariski-fpqc-quasi-coherent}{{03P2}{source}{}{sourcepointer-5}{}}
\newlabel{theorem-zariski-fpqc-quasi-coherent@cref}{{[originalreference][0][]03P2}{[1][0][]source}{}{}{}}
\makeatother
\title{Cohomology with support in a closed subscheme agrees between the Zariski and etale topologies for quasi-coherent coefficients}
\author{The Stacks Project Authors}
\date{stacks-project:etale-cohomology.tex:local-snapshot-20261011; GNU FDL 1.2 or later}
\begin{document}
\maketitle
\noindent\textit{Editorial note.} All delivered mathematical text is the author's own, copied from the pinned Stacks Project chapter file etale-cohomology.tex; no statement or proof step was rewritten, reordered or removed.\par\medskip
\noindent\textit{Editorial note.} The counted proof is the author's complete argument: the comparison of the long exact sequences for cohomology with supports, and the quasi-coherence argument for the higher direct images of the open immersion combined with the support triangle.\par\medskip
\noindent\textit{Editorial note.} The label lemma-cohomology-with-support-quasi-coherent names the counted statement, and the proof environment immediately following it in the source is delivered verbatim as its proof.\par\medskip
\noindent\textit{Editorial note.} The counted statement is carried by the author's own lemma environment in the chapter source.\par\medskip
\noindent\textit{Editorial note.} Pre-open gate: python operations/agent-authoring/tools/precard\_gate.py --source etale-cohomology.tex --label lemma-cohomology-with-support-quasi-coherent --start 14155 --end 14187 --local-id 3567 returned RESULT=OK (L3567 free, occupancy FREE, public interval clear). Candidate taken from lane-s4/halfB-priority.json.\par\medskip
\noindent\textit{Editorial note.} The statement and proof are copied from the pinned chapter file; the cited inputs are kept as external pointers with their canonical Stacks tags.\par\medskip
\noindent\textit{Allowed standard prerequisites (editorial).}\begin{itemize}
\item Agreement of Zariski, etale and fpqc cohomology for quasi-coherent modules. Same chapter, the theorem identifying the cohomology of the associated quasi-coherent sheaf on the small etale site with the usual cohomology of the module. Supplies the equalities H\textasciicircum{}q(X\_etale, F\textasciicircum{}a) = H\textasciicircum{}q(X, F) and H\textasciicircum{}q(U\_etale, F\textasciicircum{}a) = H\textasciicircum{}q(U, F) entering the long exact sequences for cohomology with supports.
\item Quasi-coherence of the higher direct images along a quasi-compact open immersion. Cohomology of Schemes, the lemma that the higher direct images of a quasi-coherent module along a quasi-compact morphism are quasi-coherent. Supplies the quasi-coherence of R\textasciicircum{}q j\_* in the second part of the lemma.
\item Higher direct images commute with the associated quasi-coherent sheaf. Descent, the lemma comparing the higher direct images of the associated quasi-coherent sheaf on the small etale site with the associated sheaf of the higher direct images. Supplies the identification of i\_*H\textasciicircum{}q\_Z(F\textasciicircum{}a) with the associated sheaf of i\_*H\textasciicircum{}q\_Z(F) in the retrocompact case.
\item The triangle relating cohomology with support to the cohomology of the open complement. Same chapter, the lemma producing the distinguished triangle computing cohomology with support in a closed subscheme. Supplies the final step translating the two identifications into the statement about the sheaves with support.
\end{itemize}\medskip
\ifdefined\hypertarget\hypertarget{sourcepointer-1}{}\fi
\noindent\textit{Source reference.} Original lemma 01XJ: sources/stacks-chow/tags.txt, source lines 2520--2520. Prerequisite pointer invoked at proof line(s) 14181; it names a standard fact of a sibling chapter that the counted argument uses, and every step of the counted proof built on it is retained verbatim in the delivered body.\par\medskip
\ifdefined\hypertarget\hypertarget{sourcepointer-2}{}\fi
\noindent\textit{Source reference.} Original lemma 071N: sources/stacks-chow/tags.txt, source lines 8997--8997. Prerequisite pointer invoked at proof line(s) 14184; it names a standard fact of a sibling chapter that the counted argument uses, and every step of the counted proof built on it is retained verbatim in the delivered body.\par\medskip
\ifdefined\hypertarget\hypertarget{sourcepointer-3}{}\fi
\noindent\textit{Source reference.} Original reference 0A45: sources/stacks-etale-cohomology/pinned-original-source.tex, source lines 14072--14072. Prerequisite pointer invoked at proof line(s) 14186; it names a standard fact of the same chapter that the counted argument uses, and every step of the counted proof built on it is retained verbatim in the delivered body.\par\medskip
\ifdefined\hypertarget\hypertarget{sourcepointer-4}{}\fi
\noindent\textit{Source reference.} Original reference 03OG: sources/stacks-etale-cohomology/pinned-original-source.tex, source lines 1332--1332. Prerequisite pointer invoked at proof line(s) 14161; it names a standard fact of the same chapter that the counted argument uses, and every step of the counted proof built on it is retained verbatim in the delivered body.\par\medskip
\ifdefined\hypertarget\hypertarget{sourcepointer-5}{}\fi
\noindent\textit{Source reference.} Original reference 03P2: sources/stacks-etale-cohomology/pinned-original-source.tex, source lines 2280--2280. Prerequisite pointer invoked at proof line(s) 14176; it names a standard fact of the same chapter that the counted argument uses, and every step of the counted proof built on it is retained verbatim in the delivered body.\par\medskip
\begin{lemma}
\label{lemma-cohomology-with-support-quasi-coherent}
Let $X$ be a scheme. Let $Z \subset X$ be a closed subscheme.
Let $\mathcal{F}$ be a quasi-coherent $\mathcal{O}_X$-module
and denote $\mathcal{F}^a$ the associated quasi-coherent sheaf
on the small \'etale site of $X$
(Proposition \ref{proposition-quasi-coherent-sheaf-fpqc}). Then
\begin{enumerate}
\item $H^q_Z(X, \mathcal{F})$ agrees with $H^q_Z(X_\etale, \mathcal{F}^a)$,
\item if the complement of $Z$ is retrocompact in $X$, then
$i_*\mathcal{H}^q_Z(\mathcal{F}^a)$ is a quasi-coherent sheaf of
$\mathcal{O}_X$-modules equal to $(i_*\mathcal{H}^q_Z(\mathcal{F}))^a$.
\end{enumerate}
\end{lemma}

\begin{proof}
Let $j : U \to X$ be the inclusion of the complement of $Z$.
The statement (1) on cohomology groups follows from the long
exact sequences for cohomology with supports and the agreements
$H^q(X_\etale, \mathcal{F}^a) = H^q(X, \mathcal{F})$ and
$H^q(U_\etale, \mathcal{F}^a) = H^q(U, \mathcal{F})$, see
Theorem \ref{theorem-zariski-fpqc-quasi-coherent}.
If $j : U \to X$ is a quasi-compact morphism, i.e., if $U \subset X$
is retrocompact, then $R^qj_*$ transforms quasi-coherent sheaves
into quasi-coherent sheaves
(Cohomology of Schemes, Lemma
\ref{coherent-lemma-quasi-coherence-higher-direct-images})
and commutes with taking associated
sheaf on \'etale sites
(Descent, Lemma \ref{descent-lemma-higher-direct-images-small-etale}).
We conclude by applying
Lemma \ref{lemma-cohomology-with-support-triangle}.
\end{proof}

\end{document}
